\heading{热学-24春-期中}
\noteinfo{题目来源于赛艇，答案由AI生成。}

\textbf{一、单选题}
\begin{question}
上传稿仅保留答案键：\textbf{C C D B D}。
\end{question}

\textbf{二、简答题}
\begin{question}
重力场中的 Boltzmann 分布及平均势能。
\begin{solution}
平衡分布
\[
f(z)=\frac{mg}{k_BT}e^{-mgz/(k_BT)},\qquad z>0.
\]
故平均重力势能
\[
\langle \varepsilon_p\rangle=\int_0^{\infty} mgz\,f(z)\,dz=k_BT.
\]
\ansmath{\langle \varepsilon_p\rangle=\int_0^{\infty} mgz\,f(z)\,dz=k_BT}
\end{solution}
\end{question}

\begin{question}
推导极端相对论性气体的压强公式。
\begin{solution}
由动量流定义压强，可写成
\[
p=n\langle mv_\alpha^2\rangle.
\]
极端相对论情形下
\[
\varepsilon_k=pc=\sum_{\alpha=x,y,z}mv_\alpha^2,
\]
各方向平均相同，因此
\[
p=\frac13 n\langle \varepsilon_k\rangle.
\]
\ansmath{p=\frac13 n\langle \varepsilon_k\rangle}
\end{solution}
\end{question}

\begin{question}
说明气球在等温大气中上升时浮力不变。
\begin{solution}
等温大气中
\[
\rho(z)=\frac{p(z)\mu}{RT}.
\]
气球内气体近似理想气体，若其内部物质的量不变，则
\[
V(z)=\frac{p_0V_0}{p(z)}.
\]
故浮力
\[
F(z)=\rho(z)gV(z)=\frac{p_0V_0\mu g}{RT},
\]
与高度无关。
\ansmath{F(z)=\rho(z)gV(z)=\frac{p_0V_0\mu g}{RT},}
\end{solution}
\end{question}

\begin{question}
泻流数与碰撞频率。
\begin{solution}
泻流数
\[
\Gamma=\frac14 n\langle u\rangle,
\qquad \langle u\rangle=\sqrt2\,\langle v\rangle.
\]
若分子有效截面为 $\sigma=\pi d^2$，则碰撞频率
\[
\langle Z\rangle=\Gamma\,4\pi d^2
=\sqrt2\,n\sigma\langle v\rangle.
\]
\ansmath{\langle Z\rangle=\Gamma\,4\pi d^2
=\sqrt2\,n\sigma\langle v\rangle}
\end{solution}
\end{question}

\begin{question}
单原子、刚性双原子、非刚性双原子的平均内能。
\begin{solution}
\[
\langle \varepsilon\rangle=
\frac32k_BT,\qquad \frac52k_BT,\qquad \frac72k_BT.
\]
（上传稿末项误写为 $7/5\,k_BT$，此处按能均分定理改正为 $7/2\,k_BT$。）
\ansmath{\langle \varepsilon\rangle=
\frac32k_BT,\qquad \frac52k_BT,\qquad \frac72k_BT}
\end{solution}
\end{question}

\begin{question}
Feynman 棘轮为什么长期平均不转？
\begin{solution}
若顺转需克服能垒 $\varepsilon$，则顺转概率与反转概率都受热涨落控制，量级同为
\[
p_+\propto e^{-\varepsilon/(k_BT)},
\qquad p_-\propto e^{-\varepsilon/(k_BT)}.
\]
因此长期平均无净转动，不能从单一热源中无代价取功。
\ansmath{p_+\propto e^{-\varepsilon/(k_BT)},
\qquad p_-\propto e^{-\varepsilon/(k_BT)}}
\end{solution}
\end{question}

\textbf{三、计算与证明题（按原稿保留结论）}
\begin{question}
题 1：由数据解得
\[
L=5^{\,t^*/100+1},
\]
并有
\[
\Delta L_1=5(5^{0.1}-1)\,\mathrm{cm},
\qquad
\Delta L_2=25(1-5^{-0.1})\,\mathrm{cm}.
\]
\end{question}

\begin{question}
题 2：由题给全微分关系
\[
\frac{\partial p}{p\,\partial T}=\frac1T,
\qquad
\frac{\partial p}{p\,\partial V}=-\frac{1}{V-b},
\]
积分得状态方程
\[
\frac{p(V-b)}{T}=C.
\]
\end{question}

\begin{question}
题 3：速率大于 $v_0$ 的分子撞击容器壁单位时间单位面积次数为
\[
\Gamma(v_0)=\frac{n}{4}\sqrt{\frac{8k_BT}{\pi m}}
\left(1+\frac{mv_0^2}{2k_BT}\right)e^{-mv_0^2/(2k_BT)}.
\]
\end{question}

\begin{question}
题 4：利用布朗分布求 Avogadro 常数，结论为
\[
N_A=\frac{3RT}{4\pi r^3(\rho-\rho_0)g\,\Delta h}\ln\frac{n_1}{n_2}.
\]
\end{question}

\begin{question}
题 5：大气温度线性递减模型 $T(z)=T_0-\alpha z$ 下，压强分布
\[
p(z)=p_0\left(1-\frac{\alpha z}{T_0}\right)^{\mu g/(R\alpha)}.
\]
当 $\alpha\to 0$ 时退化为等温结果
\[
p(z)=p_0 e^{-\mu gz/(RT_0)}.
\]
\end{question}

\begin{question}
题 6：两热容 $C_1,C_2$ 初温分别为 $\tau_{10},\tau_{20}$，传热满足 $\dot Q_{12}=\kappa(\tau_1-\tau_2)$。则温差满足
\[
\frac{d(\tau_1-\tau_2)}{dt}
=-\kappa\frac{C_1+C_2}{C_1C_2}(\tau_1-\tau_2),
\]
故半衰时间
\[
t_{1/2}=\frac{C_1C_2}{\kappa(C_1+C_2)}\ln2.
\]
此时两侧温度为
\[
\tau_{1,1/2}=\frac{(2C_1+C_2)\tau_{10}+C_2\tau_{20}}{2(C_1+C_2)},
\qquad
\tau_{2,1/2}=\frac{(C_1+2C_2)\tau_{20}+C_1\tau_{10}}{2(C_1+C_2)}.
\]
\end{question}

\begin{question}
题 7：Poiseuille 定律。
\begin{solution}
由牛顿黏滞定律得速度分布
\[
u(\rho)=\frac{\Delta p}{4\eta L}(r^2-\rho^2),
\]
积分得体积流量
\[
Q=\int_0^r 2\pi\rho\,u(\rho)\,d\rho
=\frac{\Delta p\,\pi r^4}{8\eta L}.
\]
\ansmath{Q=\int_0^r 2\pi\rho\,u(\rho)\,d\rho
=\frac{\Delta p\,\pi r^4}{8\eta L}}
\end{solution}
\end{question}


